% Options for packages loaded elsewhere
\PassOptionsToPackage{unicode}{hyperref}
\PassOptionsToPackage{hyphens}{url}
%
\documentclass[
]{article}
\usepackage{amsmath,amssymb}
\usepackage{iftex}
\ifPDFTeX
  \usepackage[T1]{fontenc}
  \usepackage[utf8]{inputenc}
  \usepackage{textcomp} % provide euro and other symbols
\else % if luatex or xetex
  \usepackage{unicode-math} % this also loads fontspec
  \defaultfontfeatures{Scale=MatchLowercase}
  \defaultfontfeatures[\rmfamily]{Ligatures=TeX,Scale=1}
\fi
\usepackage{lmodern}
\ifPDFTeX\else
  % xetex/luatex font selection
\fi
% Use upquote if available, for straight quotes in verbatim environments
\IfFileExists{upquote.sty}{\usepackage{upquote}}{}
\IfFileExists{microtype.sty}{% use microtype if available
  \usepackage[]{microtype}
  \UseMicrotypeSet[protrusion]{basicmath} % disable protrusion for tt fonts
}{}
\makeatletter
\@ifundefined{KOMAClassName}{% if non-KOMA class
  \IfFileExists{parskip.sty}{%
    \usepackage{parskip}
  }{% else
    \setlength{\parindent}{0pt}
    \setlength{\parskip}{6pt plus 2pt minus 1pt}}
}{% if KOMA class
  \KOMAoptions{parskip=half}}
\makeatother
\usepackage{xcolor}
\setlength{\emergencystretch}{3em} % prevent overfull lines
\providecommand{\tightlist}{%
  \setlength{\itemsep}{0pt}\setlength{\parskip}{0pt}}
\setcounter{secnumdepth}{-\maxdimen} % remove section numbering
\ifLuaTeX
  \usepackage{selnolig}  % disable illegal ligatures
\fi
\IfFileExists{bookmark.sty}{\usepackage{bookmark}}{\usepackage{hyperref}}
\IfFileExists{xurl.sty}{\usepackage{xurl}}{} % add URL line breaks if available
\urlstyle{same}
\hypersetup{
  hidelinks,
  pdfcreator={LaTeX via pandoc}}

\author{}
\date{}

\begin{document}

\hypertarget{untitled}{%
\section{Untitled}\label{untitled}}

Come vediamo? L'occhio è una fotocamera avanzata
Nell'occhio, la luce entra, viene messa a fuoco da una lente (il cristallino) e viene infine proiettata su un sensore (la retina)

La luce entra nel ``pinhole'' e sulla parete opposta si forma l'immagine capovolta:

\begin{figure}
\centering
{image.png}
\caption{image.png}
\end{figure}

\begin{figure}
\centering
{image.png}
\caption{image.png}
\end{figure}

I raggi vanno dal basso verso l'alto e dall'alto verso il basso

Nel pinhole:

\begin{itemize}
\tightlist
\item
  Per avere una immagine nitida serve un foro piccolo ma l'immagine è quasi un punto
\item
  Se invece abbiamo un foro grande entra tanta luce e l'immagine appare come una macchia sfocata. Questo fa riferimento al fenomeno di diffrazione.
\end{itemize}

Per risolvere il problema della legge di sovrapposizione si fa riferimento alla lente sottile

\begin{enumerate}
\def\labelenumi{\arabic{enumi}.}
\tightlist
\item
  Un apertura grande raccoglie la luce
\item
  Tramite rifrazione fa coincidere tutti i punti in un unico punto nitido detto fuoco
\end{enumerate}

Quindi piega tutti raggi e li fa coincidere in un unico punto

\[
\frac{1}{u} + \frac{1}{v} = \frac{1}{f}
\]

U è la distanza tra la lente e l'oggetto

v è la distanza tra la l ente e il punto dove si forma l'immagine

f è la lunghezza focale

\end{document}
